%% ==========================================================================
\section{The obscurity class, and what vendors claim for it}
\label{sec:obscurity}
%% ==========================================================================


The designs above interpose a delay or a veto. A large and widely shipped class
does something categorically different: it \emph{erases}. An alternate PIN, when
entered, destroys the device's copy of the secret. Blockstream's Jade documents
such a PIN under both of the names its help center uses for it --- \emph{duress}
and \emph{wallet-erase} --- as an alternate PIN that ``automatically delete[s]
the encrypted wallet information from Jade if entered,'' presented as
``especially useful in the event of a physical threat, as you can provide this
PIN to an attacker.'' On entry the device reports only ``Internal Error,'' so
that the erasure reads as a
malfunction~\cite{jade-duress-pin,jade-erase-pin}.

Grant the design the strength it genuinely has, and which the delay
architectures above lack. The erasure is \emph{instantaneous and covert}:
$\Attacker$ learns of it only once it is complete, so there is no interval in
which he can intervene to prevent it. Nothing in what follows turns on
$\Attacker$ out-racing the mechanism, because he cannot.

\medskip\noindent\textbf{The claim decomposes, and only one half is real.} What
the design offers $\Victim$ is not one mechanism but two, stacked and easily
conflated:
\begin{enumerate}[nosep,leftmargin=2.4em,label=(\alph*),ref=\alph*]
  \item\label{sd:erase} \emph{the erasure} --- a genuine, irreversible state
        change, which happens whatever $\Attacker$ knows or believes; and
  \item\label{sd:belief} \emph{the belief} that nothing further remains to be
        had --- manufactured by the ``Internal Error'' report and the return to
        setup, which are engineered to read as breakage rather than as a move.
        This is a \emph{presentation} choice and not a property of erasure as
        such; Remark~\ref{rem:vendor-claims} records a vendor asked directly for
        it who declined.
\end{enumerate}
Only~\ref{sd:erase} is Kerckhoffs-independent. The whole question is how much
coercion resistance it carries on its own, and the answer, from the vendor's own
instructions, is \emph{none}.

The same documentation warns that afterwards ``there is no way to regain access
to your funds except by restoring using your recovery phrase,'' and directs the
holder to ``make sure you have a backup before
proceeding''~\cite{jade-erase-pin}. The feature therefore \emph{presupposes} a
surviving backup: a holder who follows the documentation has one, and a holder
who does not has destroyed the stash rather than defended it. What
erasure removes is thus not the secret but the \emph{device path} to it --- and a
path that cost $\Attacker$ nothing to lose, since \ref{sd:erase} leaves
$\Victim$ alive, present, and still holding everything needed to restore. In the
ideal model \ref{sd:erase} contributes exactly nothing to the floor, for the same
reason a stage below the memory-hard chain contributes nothing to it
(the shallow-stage null result in~\cite{deadly-race}): what is reachable without work is not a barrier.

All of the design's protective content therefore sits in~\ref{sd:belief}, which
is an obscurity bet and nothing else. That places the design squarely in the
naive tier (Definition~\ref{def:naive-cr}): what protection it offers is
contingent on $\Attacker$'s ignorance of a documented feature, which is to say it
forfeits non-obscurity (Property~\ref{p:nonob}). It is worth being exact about
which failure this is. The design does \emph{not} fail the No-Gray-Area
Criterion --- an erasure destroys a path and names no second racer, so
Lemma~\ref{lem:drl} does not reach it, and its row in the comparison is DR-safe
and condemned regardless.

NNPP (Assumption~\ref{lem:nnpl}) is what makes the obscurity bet unwinnable:
$\Attacker$ cannot verify that no further copy exists, and here he need not even
reason his way to the suspicion, since the manual instructs the holder to keep a
backup. He is therefore correct, by the design's own construction, to read the
erasure as the closing of one route rather than the end of the matter. The route
that remains runs through $\Victim$'s memory of a seed phrase, or of where a
backup is kept: the erasure deletes the path that requires hurting no one and
retains the path that does.

\begin{remark}[The obscure self-destruct is the obscure scheme with no winning branch]
\label{rem:two-ignorances}
Obscure defenses are usually compared by \emph{how likely} $\Attacker$ is to be
ignorant. The sharper comparison is \emph{how many} independent ignorances each
one needs, and here the self-destruct is strictly the worse bet --- worse, in
particular, than the decoy it is commonly sold beside.

A decoy needs one premise: that $\Attacker$ does not know decoys exist. Grant him
that ignorance and the decoy \emph{terminates the encounter} --- he leaves with a
plausible balance and a story that closes.

A self-destruct needs two, and they are independent. It needs $\Attacker$ ignorant
of the feature (sub-Kerckhoffs), so that he reads the failure as breakage rather
than as a deliberate move; and it needs him ignorant of NNPP's consequence
(sub-NNPP), so that he does not reason to an undisclosed backup. The second is
not implied by the first, and this is the design's central weakness: \emph{even
total success at the first leaves the encounter open.} ``This device is broken''
does not entail ``this person has no bitcoin''; it entails that the instrument
for reaching it is broken, whose natural sequel is the question \emph{where is
your backup?} --- addressed to a $\Victim$ who is still present and still knows.

So the self-destruct fails against an informed $\Attacker$ for the ordinary
Kerckhoffs reason, and against an ignorant one it fails anyway, because the
belief it manufactures does not discharge the encounter. It is the one obscure
scheme with no branch on which it wins.

We did not have to derive that regress: a vendor documents it. Coldcard ships
trick PINs whose actions include a duress wallet --- ``a personal safety
feature\ldots{} If you must reveal a PIN under duress, give the duress PIN
instead of your Main PIN'' --- alongside \emph{Brick Self} and \emph{Wipe
Seed}~\cite{coldcard-trick-pins}. Its firmware documentation then addresses the
case where the decoy is not believed, and its advice is to escalate to
destruction: ``you may be asked what the actual duress PIN is, while under
duress. We suggest providing the `brickme' PIN in that
case''~\cite{coldcard-pin-entry}. That is the manufacturer contemplating an
$\Attacker$ who knows decoys exist, and recommending, as the answer, an
irreversible erasure performed with $\Victim$ still in the room and still able
to be asked for a backup --- the ignorance-count argument above, arrived at
independently and published as guidance. The same document instructs the holder
to ``triple-check your backups are correct'' before enabling that PIN, so the
backup presupposition of \S\ref{sec:obscurity} is not peculiar to one vendor
either.

The same page goes further, in two admissions worth quoting because a vendor
rarely makes them. The first concedes that the decoy is distinguishable to the
$\Attacker$ who bothers: ``if you are somehow facing an attacker who is willing
to verify he has the real main PIN, it's possible that careful analysis of system
responses will imply he's working with the duress PIN.'' The remedy offered is
not a repair but a request --- ``if you discover any sequences that reveal this
easily, please tell us and we'll see if we can cover them up
better''~\cite{coldcard-pin-entry}. That is obscurity stated as a maintenance
programme, and it is the clearest statement we have found of what
Property~\ref{p:nonob} rules out. What is being defended is not that the decoy is
indistinguishable, but that no one has yet published how to distinguish it; the
defence is administered by patching the disclosures as they arrive. A scheme
whose security is the current state of the disclosure backlog is exactly the
scheme whose guarantee $\Victim$ cannot invoke in the room, because $\Victim$
cannot know which sequences $\Attacker$ has read.

The second admission is the denial move itself, offered where bricking is
unpalatable: ``alternatively, you can say you set the duress PIN once, but have
since forgotten it.'' The escalation to destruction and this are the document's
two exits, and both are the naive tier. The claim \emph{I have forgotten it} is a
proposition about $\Victim$'s own knowledge, and NNPP
(Assumption~\ref{lem:nnpl}) is precisely the observation that $\Victim$ cannot
close it: there is no move that converts a true \emph{I cannot} into an
$\Attacker$-verifiable one, and by Definition~\ref{def:naive-cr} the scheme's
fate on that branch is decided by whether he believes it. The vendor has
enumerated the branches and found, on each, a bluff.

That the line is published guidance sharpens rather than softens it, in the way
Remark~\ref{rem:decoy-market} describes for the coached decoy. A script printed
in a manufacturer's manual is a script $\Attacker$ can read as easily as
$\Victim$; the sentence that teaches the excuse also teaches that the excuse is
taught. Its credibility declines in proportion to its circulation, and the
holder who most needs it inherits the discount without ever having been told it
exists.
\end{remark}

\begin{remark}[The class, and what each vendor actually claims for it]
\label{rem:vendor-claims}
Three vendors ship the pair, which is what makes this a design consensus rather
than one firm's misstep; but they do not make the same claims, and the
difference is worth recording precisely. Blockstream markets the \emph{erasure}
for physical threat while claiming nothing of the kind for its passphrase
wallets. Coldcard markets \emph{both} halves, describing the duress wallet as a
safety feature and the brick PIN as the fallback when it fails. Trezor is the
mirror image of Blockstream: its passphrase wallets are sold on deniability,
while its wipe code --- which the vendor itself calls a ``self-destruct'' PIN,
and which erases the recovery seed with it --- is documented without any mention
of duress, coercion or physical threat~\cite{trezor-wipe-code}. That vendor was
moreover asked, in January 2022, to make the wipe \emph{concealable}: a
user-defined message on entry, so that an erasure could be presented as a
hardware fault rather than announcing itself. The request was closed as \emph{not
planned}, and the device still reports the wipe
plainly~\cite{trezor-wipe-deniability}. Component~\ref{sd:belief} is therefore
not intrinsic to the self-destruct class at all: it is a presentation one vendor
builds and another, asked for it directly, refused. What \S\ref{sec:obscurity}
analyses under~\ref{sd:belief} is specific to the ``Internal Error'' design, and
a wipe that announces itself carries \ref{sd:erase} alone --- which, by the
argument above, is no coercion resistance at all, and is not sold as any. Two of the three
therefore decline to make the coercion claim on one of their two mechanisms. We
note it because the criterion is not advanced against vendors, and a vendor that
ships a mechanism without claiming it defeats coercion has not made the error
this section is about.
\end{remark}

\begin{remark}[Two bluffs the erasure invites]\label{rem:erase-bluffs}
An $\Attacker$ who knows the design (Assumption~\ref{ass:power}) need not walk
into the erasure at all. Two alternatives stand open to him, each a wager rather
than a certainty:
\begin{enumerate}[nosep,leftmargin=1.6em]
  \item \emph{Take the device hostage up front.} Rather than enter any PIN,
        seize the device and hold it as a hostage token (the hostage-token property in~\cite{deadly-race},
        the hostage-token remark in~\cite{deadly-race}). An erase PIN is irrelevant to a coercer who
        never authenticates; and by the backup-probe remark in~\cite{deadly-race} the terms on
        which $\Victim$ ransoms the device are themselves a probe for the very
        backups $\Victim$ cannot disprove.
  \item \emph{Coerce against deployment.} Brutalize $\Victim$ to deter the PIN's
        use in the first place, wagering that a holder under sufficient duress
        will not spend the one irreversible move available to them.
\end{enumerate}
Neither wager is sure to pay, and that is precisely the point: both are bluff
games, and the holder who has bet on an erasure has entered one without holding
the means to end it. A $\Victim$ who erases and then \emph{truthfully} denies
retaining any backup arrives at the position the registry already records --- the
streamer whose denial was simply disbelieved, and who was tortured for hours
regardless~\cite{ratirl} (\S\ref{sec:registry}).
\end{remark}

The ``Internal Error'' disguise inherits the same objection from a different
direction: it is a bet on $\Attacker$'s ignorance of a consumer product, and it
degrades exactly as that product's own public documentation becomes known
(the Four-Properties development in~\cite{deadly-race}). The class therefore lands in the \textbf{naive} tier of
the time-gating tiers in~\cite{deadly-race} precisely as that definition describes it --- an
obscure self-destruct, safe only against a sub-Kerckhoffs attacker --- and it is
worth keeping the self-destruct distinct from the \emph{cancel} key named
alongside it there. A cancel reverses a pending action and thereby names a second
racer; a self-destruct destroys a path and names none. Only the former is subject
to Lemma~\ref{lem:drl}, which is why the two fail differently despite sharing a
tier. the perceptual-oracle variants in~\cite{deadly-race} has already disposed of the tier as a
whole, by the dilemma that an $\Attacker$ whose ignorance can be assumed makes
the elaborate mechanism unnecessary, and one whose ignorance cannot makes it
ineffective; the present subsection is that dilemma instantiated in a shipping
product.

\medskip\noindent\textbf{The companion decoy, classified separately.} The same
device ships a second and quite different obscure mechanism, and the two should
not be run together. A BIP-39 passphrase ``is added to your Jade's
recovery phrase or SeedQR and is used to create your wallet''; ``any number of
passphrases can be entered to change the wallet you will log into''; and the
device ``will not store your passphrase''~\cite{jade-passphrase}. Distinct
passphrases therefore yield distinct wallets from one recovery phrase, and the
familiar community practice is to keep a low-balance wallet at the empty
passphrase to be surrendered under coercion.

Two observations. First, and to the vendor's credit, \emph{Blockstream does not
market it this way}: the passphrase documentation claims no deniability, decoy
or duress use, in pointed contrast to the erase-PIN pages, which name physical
threat explicitly. The duress framing here is the community's,
not the manufacturer's, and the distinction is worth preserving in any account of
who claimed what.

Second, the mechanism is the plausible-deniability layer already treated in
the Four-Properties development in~\cite{deadly-race} and dismissed for the same Kerckhoffs reason: it rests on
$\Attacker$'s ignorance, and an $\Attacker$ who knows that passphrase wallets
exist simply persists past the surrendered one, with NNPP guaranteeing $\Victim$
cannot close the question. It is not, however, the \emph{same} failure as the
self-destruct's, and by Remark~\ref{rem:two-ignorances} it is the less bad of the
two: the decoy at least has a branch on which it wins, since an $\Attacker$ who
believes it has been paid and can leave. The self-destruct has none. That a
single product ships both, marketing only the weaker one for physical threat, is
the clearest illustration in this section of how little the design space has been
mapped against the Criterion.

\begin{remark}[A vendor-documented instance of the naive tier]\label{rem:hidden-medium}
The decoy's duress framing was the community's, not the manufacturer's. One vendor
does make the claim in its own documentation. Sparrow's FAQ, answering ``How can I
choose where my wallet is saved?'', documents a data-directory flag and adds that
``[t]his feature allows you to store all Sparrow data on removable media making for
more plausible deniability''~\cite{sparrow-faq}.

The feature does not guard the asset. Sparrow's own best-practice page has the seed
words backed up separately from the device, ``ideally in a different
location''~\cite{sparrow-best-practices} --- so the concealed directory is neither
necessary nor sufficient for $\Kset$. A Kerckhoffs $\Attacker$ coerces the seed
backup, loads it into a wallet of his own, and never touches the concealed medium at
all. What is left for the medium to do is carry the holder's belief that, having
surrendered what was found, they may now say there is nothing more. The decoy at
least stands between $\Attacker$ and the coins; this stands nowhere near them.

That is Definition~\ref{def:naive-cr} exactly --- a guarantee contingent on
$\Attacker$'s ignorance, here of a \emph{documented} feature --- and the continuation
is the one the registry records. A denial offered into the prevailing norm is
pre-discredited (Corollary~\ref{cor:disclosure}), so the rational continuation of an
interrogation that opens with one is more coercion, and the truthful denier is
brutalized alongside the lying one (\S\ref{sec:registry}). The feature adds nothing
but the confidence to make the denial.
\end{remark}

\begin{remark}[The decoy is sold, and coached for credibility]\label{rem:decoy-market}
The decoy is not merely a practice the community tolerates; it is a product and a
curriculum. A Spanish-language self-custody consultancy lists, among paid service
tiers, ``\emph{Passphrase (tu b\'oveda se\~nuelo)} --- Billetera se\~nuelo +
billetera real. Protecci\'on ante coacci\'on/coerci\'on'' --- a decoy wallet
marketed, in as many words, as protection against coercion~\cite{tlt-servicios}.
Published guidance goes further and coaches the performance: one guide has the
holder keep a small balance on the bare seed as ``what you'd hand over under
physical coercion'', and adds that ``for this to work, the decoy balance must be
\emph{believable} --- enough to look like a real wallet, not enough to devastate
you if lost''~\cite{bsg-passphrase}. The counsel is also delivered live: one of us
appeared as guest on a self-custody programme~\cite{ong-autocustodia} in which the
host worked through a decoy-construction walkthrough over the guest's contrary
advice.

That last instruction is the naive tier stated as a training objective. A decoy
need be believable only to an $\Attacker$ who might disbelieve it, and the
counsel is silent on what happens when he does --- which, by
Definition~\ref{def:naive-cr}, is the whole question. A better-rehearsed lie
raises the chance $\Attacker$ takes the bait at the margin; it does not create a
branch on which the scheme wins against an $\Attacker$ who knows decoys exist,
and NNPP (Assumption~\ref{lem:nnpl}) guarantees $\Victim$ cannot close that
question however well the performance goes. What is sold as a mechanism is a
wager on the attacker's ignorance, priced and taught as though it were a control.

The counsel is also self-defeating in aggregate. Publicised training raises the
\emph{prior} that a fluent denial was rehearsed: once $\Attacker$ knows credibility
is taught and sold, fluency stops being evidence of truth and becomes evidence of
coaching. The remedy degrades the thing it sells --- and not only for its
customers. A holder who never trained now denies into an $\Attacker$ who discounts
fluent denial generally, so the cost falls on people who bought nothing. This is
the Obscurity Paradox (Remark~\ref{rem:obscurity-paradox}) run backwards: there, a
public non-obscure design makes $\Victim$'s claims \emph{more} plausible; here, a
public obscure training makes them \emph{less}.
\end{remark}

\begin{remark}[The Obscurity Paradox --- non-obscure designs strengthen obscure ones]
\label{rem:obscurity-paradox}
TGPO enters neither bluff. There is nothing to erase, the deep secret having
never been on the device (the floor theorem in~\cite{deadly-race}), and so nothing whose
non-deployment can be coerced. But the comparison yields more than the absence
of a failure mode, and the surplus runs in a direction the obscurity argument
would not predict.

An obscure defense asks $\Attacker$ to believe a \emph{claim} --- that this is
the only wallet, that nothing further remains --- against both an incentive to
disbelieve it and, by NNPP, no means of settling it. A non-obscure defense
changes what making that claim costs, and it does so by the mechanism set out in
the adversary-education remark in~\cite{deadly-race}: an $\Attacker$ who credits that an
effective mechanism is in place is crediting that lying is not $\Victim$'s only
available tool, and a holder with effective alternatives has correspondingly
less reason to resort to one.

The ordering this inverts is worth stating plainly, and we name it \textbf{the
Obscurity Paradox}: \emph{a publicly specified, Kerckhoffs-clean design improves
the standing of precisely those obscure denials it was built to make
unnecessary.} Obscurity remains unsound as a primary defense, for the reasons
the Four-Properties development in~\cite{deadly-race} gives. It is not, however, indifferent which design it
accompanies --- a denial made from within a mechanism that removes the need to
lie is a stronger denial than the same words spoken from within one that
requires it. The same asymmetry has an epistemic cost, taken up next
(Remark~\ref{rem:anosognosia}).
\end{remark}
